%%%%%%%%%%%%%%%%%%%%%%%%%%%%%%%%%%%%%%%%%%%%%%%%%%%%%%%%%%%%%%%%%%%%%%%%%%%%%%%
%% Inputs
%%%%%%%%%%%%%%%%%%%%%%%%%%%%%%%%%%%%%%%%%%%%%%%%%%%%%%%%%%%%%%%%%%%%%%%%%%%%%%%

\documentclass[
ngerman          			% neue deutsche Rechtschreibung
,a4paper          			% Papiergrösse
%,twoside          			% Zweiseitiger Druck (rechts/links)
%,10pt             			% Schriftgrösse
%,11pt
, 12pt
,pdftex
%,disable         			% Todo-Markierungen auschalten
]{article}
\usepackage[utf8]{inputenc}
\usepackage[ngerman]{babel}
\usepackage[T1]{fontenc}
\usepackage{amsmath}
\usepackage{amsfonts}
%\usepackage{subfigure}
\usepackage{framed}
\usepackage{booktabs}
\usepackage{setspace}
\usepackage{threeparttable}
\usepackage{amssymb}
\usepackage{graphicx}
\usepackage[headings]{fullpage}
\usepackage{fancyhdr}
\usepackage[hyphens]{url}
\usepackage[toc, acronym]{glossaries}
\usepackage[
	hyperref=true,			% Klickbare Referenzen in der PDF-Datei
	backref=true,			% In der Literaturref. die Seiten angeben, wo ein \cite dazu steht
	bibencoding=inputenc,	% s. inputenc-Paket
	style=ieee,
	backend=bibtex,
	sorting=nty]{biblatex}
%\addbibresource{Bericht.bib}
\let\cite\parencite

\usepackage{icomma}
\usepackage{float}
\usepackage{pdfpages}
\usepackage{hyperref}
\usepackage{csquotes}
\usepackage{listings}
\usepackage[left=2.5cm,right=2.5cm,top=2cm,bottom=3.5cm]{geometry}
\onehalfspacing 			% Zeilenabstand


%%%%%%%%%%%%%%%%%%%%%%%%%%%%%%%%%%%%%%%%%%%%%%%%%%%%%%%%%%%%%%%%%%%%%%%%%%%%%%%
%% Angaben zur Arbeit
%%%%%%%%%%%%%%%%%%%%%%%%%%%%%%%%%%%%%%%%%%%%%%%%%%%%%%%%%%%%%%%%%%%%%%%%%%%%%%%

\newcommand{\Autor}{Max Musternann}
\newcommand{\MatrikelNummer}{XXXXXXX}
\newcommand{\Kursbezeichnung}{TINF1XBX}

\newcommand{\FirmenName}{Musterfirma}
\newcommand{\FirmenStadt}{Musterort}
\newcommand{\FirmenLogoDeckblatt}{\includegraphics[width=2cm]{logo}}

\newcommand{\BetreuerFirma}{Moritz Mustermann}
\newcommand{\BetreuerDHBW}{Dodge}

\newcommand{\Was}{Bachelor-/Studien-/Projektarbeit}

\newcommand{\Titel}{Hier steht ein Mustertitel meiner Bachelor-/Studien-/Projektarbeit}
\newcommand{\AbgabeDatum}{xx. November 109 B.C.}

\newcommand{\Dauer}{von wann bis wann}
\newcommand{\Abschluss}{Bachelor of Science}

\newcommand{\Studiengang}{Informatik / Informationstechnik}

\hypersetup{%%
	pdfauthor={\Autor},
	pdftitle={\Titel},
	pdfsubject={\Was}
}


\makeglossaries
\newacronym{ros}{ROS}{Robot Operating System}

\newglossaryentry{report}
{
    name=Bericht,
    description={Eine strukturierte Darstellung von Inhalten aus Datenbeständen}
}

%%%%%%%%%%%%%%%%%%%%%%%%%%%%%%%%%%%%%%%%%%%%%%%%%%%%%%%%%%%%%%%%%%%%%%%%%%%%%%%

\bibliography{Bericht}

%%%%%%%%%%%%%%%%%%%%%%%%%%%%%%%%%%%%%%%%%%%%%%%%%%%%%%%%%%%%%%%%%%%%%%%%%%%%%%%
%% Dokument
%%%%%%%%%%%%%%%%%%%%%%%%%%%%%%%%%%%%%%%%%%%%%%%%%%%%%%%%%%%%%%%%%%%%%%%%%%%%%%%


\begin{document}
	
	
	\begin{titlepage}
		\begin{center}
			\vspace*{-2cm}
			\FirmenLogoDeckblatt\hfill\includegraphics[width=4cm]{dhbw-logo}\\[2cm]
			{\Huge \Titel}\\[1cm]
			{\Huge\scshape \Was}\\[1cm]
			{\large für die Prüfung zum}\\[0.5cm]
			{\Large \Abschluss}\\[0.5cm]
			{\large des Studienganges \Studiengang}\\[0.5cm]
			{\large an der}\\[0.5cm]
			{\large Dualen Hochschule Baden-Württemberg Karlsruhe}\\[0.5cm]
			{\large von}\\[0.5cm]
			{\large\bfseries \Autor}\\[1cm]
			{\large Abgabedatum \AbgabeDatum}
			\vfill
		\end{center}
		\begin{tabular}{l@{\hspace{2cm}}l}
			Bearbeitungszeitraum			& \Dauer			\\
			Matrikelnummer					& \MatrikelNummer	\\
			Kurs							& \Kursbezeichnung	\\
			Ausbildungsfirma				& \FirmenName		\\
											& \FirmenStadt		\\
			Betreuer der Ausbildungsfirma	& \BetreuerFirma	\\
			Gutachter der Studienakademie	& \BetreuerDHBW		\\
		\end{tabular}
	\end{titlepage}

	%%%%%%%%%%%%%%%%%%%%%%%%%%%%%%%%%%%%%%%%%%%%%%%%%%%%%%%%%%%%%%%%%%%%%%%%%%%%%%%
	%%%%%%%% Erklärung %%%%%%%%
	
	%%%%%%%%%%%%%%%%%%%%%%%%%%%%%%%%%%%%%%%%%%%%%%%%%%%%%%%%%%%%%%%%%%%%%%%%%%%%%%%
%% Descr:       Vorlage für Berichte der DHBW-Karlsruhe, Erklärung
%% Author:      Prof. Dr. Jürgen Vollmer, vollmer@dhbw-karlsruhe.de
%% $Id: erklaerung.tex,v 1.11 2020/03/13 14:24:42 vollmer Exp $
%% -*- coding: utf-8 -*-
%%%%%%%%%%%%%%%%%%%%%%%%%%%%%%%%%%%%%%%%%%%%%%%%%%%%%%%%%%%%%%%%%%%%%%%%%%%%%%%

% In Bachelorarbeiten muss eine schriftliche Erklärung abgegeben werden.
% Hierin bestätigen die Studierenden, dass die Bachelorarbeit, etc.
% selbständig verfasst und sämtliche Quellen und Hilfsmittel angegeben sind. Diese Erklärung
% bildet das zweite Blatt der Arbeit. Der Text dieser Erklärung muss auf einer separaten Seite
% wie unten angegeben lauten.

\newpage
\thispagestyle{empty}
\begin{framed}
\begin{center}
\Large\bfseries Erklärung
\end{center}
\medskip
\noindent
% siehe §5(3) der \enquote{Studien- und Prüfungsordnung DHBW Technik} vom 29.\,9.\,2017 und Anhang 1.1.13
Ich versichere hiermit, dass ich meine \Was mit dem Thema:
\enquote{\Titel}
selbstständig verfasst und keine anderen als die angegebenen Quellen und Hilfsmittel benutzt habe. Ich versichere zudem, dass die eingereichte elektronische Fassung mit der gedruckten Fassung übereinstimmt.
\vspace{3cm}
\noindent \\
\underline{\hspace{4cm}}\hfill\underline{\hspace{6cm}}\\
Ort~~~~~Datum\hfill Unterschrift\hspace{3.8cm}
\end{framed}

\vfill
%\emph{Sofern  vom Dualen Partner ein Sperrvermerk gewünscht wird, ist folgende Formulierung zu verwenden:}
\begin{framed}
\begin{center}
\Large\bfseries Sperrvermerk
\end{center}
\medskip
\noindent
Der Inhalt dieser Arbeit darf weder als Ganzes noch in Auszügen Personen
außerhalb des Prüfungsprozesses und des Evaluationsverfahrens zugänglich gemacht
werden, sofern keine anderslautende Genehmigung vom Dualen Partner vorliegt.
\end{framed}

%%%%%%%%%%%%%%%%%%%%%%%%%%%%%%%%%%%%%%%%%%%%%%%%%%%%%%%%%%%%%%%%%%%%%%%%%%%%%%%
\endinput
%%%%%%%%%%%%%%%%%%%%%%%%%%%%%%%%%%%%%%%%%%%%%%%%%%%%%%%%%%%%%%%%%%%%%%%%%%%%%%%

	
	%%%%%%%% Verzeichnisse %%%%%%%%
	
	\pagenumbering{Roman}
	
	\clearpage
	\tableofcontents			% Inhaltsverzeichnis hier ausgeben
	\clearpage
	\listoffigures				% Liste der Abbildungen
	\clearpage
	\listoftables
	%\clearpage					% Liste der Tabellen
	%\listofequations			% Liste der Formeln
	\clearpage
	
	%%%%%%%% Abstract %%%%%%%%
	\selectlanguage{ngerman}
	\begin{abstract}
		\centering Stand \today
		\thispagestyle{empty}
	\end{abstract}
	\clearpage
%	\selectlanguage{english}
%	\begin{abstract}
%		\centering State \today
%		\thispagestyle{empty}
%	\end{abstract}
%	\clearpage
%	\selectlanguage{ngerman}
	
	%%%%%%%% Styling %%%%%%%%
	
	\pagenumbering{arabic}
	\setcounter{page}{1}
	
	%\pagestyle{fancy}
	%\fancyhf{}
	%\setlength{\headheight}{35pt}
	
	\cfoot{\thepage}
	
	%%%%%%%% Kapitel %%%%%%%%
	
	\clearpage
	\section{Einleitung}
	\label{sec:Einleitung}
	Hier steht meine Einleitung \cite{template.2021}
	
	\clearpage
	\section{Grundlagen}
	\label{sec:Grundlagen}
	%\input{Grundlagen/Grundlagen.tex}
	
	\clearpage
	\section{Umsetzung}
	\label{sec:Umsetzung}
	%\input{Umsetzung/Umsetzung.tex}
	
	\clearpage
	\section{Ausblick}
	\label{sec:Ausblick}
	%\input{Ausblick/Ausblick.tex}
	
	%%%%%%%% Anhang %%%%%%%%
	
	\clearpage
	\pagenumbering{Roman}
	\setcounter{page}{8}	% passende Zahl einsetzen
	
	\appendix
	\section{test}
	
	%%%%%%%% Literaturverzeichnis %%%%%%%%
	\clearpage
	\addcontentsline{toc}{section}{Literaturverzeichnis}
	
	\def\refname{Literaturverzeichnis}
	\printbibliography
	%%%%%%%% Abkürzungen und Glosar %%%%%%%%
	\printglossary[title={Abkürzungsverzeichnis},type=\acronymtype]
\end{document}

%%%%%%%%%%%%%%%%%%%% Frequently used commands %%%%%%%%%%%%%%%%%%%%


%\noindent
%\hangindent1cm
%\textbf{Keyword} und der dazugehörige, eingerückte Text


%\begin{figure}[tbt]
%	\begin{center}
%		\includegraphics[width=8cm]{images/}
%	\end{center}
%	\caption[Das steht im Verzeichnis]{Das steht unterm Bild (Abbildung aus \cite{Quelle})}
%	\label{fig:figure}
%\end{figure}